\chapter*{Kurzfassung}
\thispagestyle{empty}
Diese Bachelorarbeit untersucht die Entwicklung und Leistungsfähigkeit 
moderner Chatbots unter Verwendung verschiedener Large Language Models. 
Der Schwerpunkt liegt auf der Konzeption und der prototypischen 
Implementierung eines Chatbots, sowie dem anschließenden Vergleich von LLMs
innerhalb des entwickelten Systems. \par\par
\bigskip
Die Konzeption zeigte die Vielzahl aktueller 
Möglichkeiten zur Chatbot-Entwicklung. Die Implementierung 
wurde mit dem Framework Langchain durchgeführt, 
wodurch die Benutzerfreundlichkeit moderner Entwicklungstools 
deutlich wurde. Der resultierende Chatbot wurde 
erfolgreich in Google Chat integriert. \par
\bigskip
Im Anschluss wurden vier führende LLMs 
(Llama3, Phi-3, Mistral und Gemma 2) verglichen. 
Der Vergleich erfolgte durch eine quantitative 
Analyse mittels des Frameworks RAGAS, sowie einer
qualitative Bewertung durch Fragebögen. 
In beiden Vergleichsansätzen erwies sich Llama3 
als das beste LLM. \par
\bigskip
Diese Arbeit liefert wertvolle Einblicke in die praktische 
Umsetzung und Leistungsfähigkeit moderner Chatbot-Technologien 
und bietet eine Grundlage für zukünftige Entwicklungen in diesem Bereich.

\bigskip

\noindent
Schlagworte: Chatbot, Generative KI, Large-Language-Modell, Prototyp, 
Retrieval-Augmented-Generation, Vektordatenbanken

